\section{Future DESTINO PoC Integration Blueprint}
\label{sec:detailed_design_blueprint}

This section records the intended integration shape for the future proof-of-concept without choosing the PoC metric.
Its purpose is to keep Phase I connected to the deployment target while preserving the decision rule: empirical comparison and
feasibility scoring must happen before implementation.

\subsection{Design Objective and Boundary}
\label{sec:design_objective_boundary}

The future PoC should support evidence-based developer reputation records that are traceable to hash-pinned software artifacts.
The contract layer should not compute repository metrics directly. Repository traversal, Git history analysis, duration
normalization, bot handling, and ownership calculations remain off-chain because they require file-system access, Git tooling,
and potentially unbounded iteration. The contract stores the result and the commitments needed to audit it.

The boundary is:
\begin{itemize}[leftmargin=*]
  \item \textbf{Off-chain:} artifact retrieval, hash verification, metric extraction, cleaning, normalization, and report generation;
  \item \textbf{On-chain:} submission registration, compact metric-record storage, evidence-hash storage, permission checks, and events;
  \item \textbf{Not in Phase I:} final metric selection, contract implementation, and M5 maintenance-responsiveness implementation.
\end{itemize}

\subsection{Metric-Agnostic Record Shape}
\label{sec:design_record_shape}

The future contract interface can remain metric-agnostic until Phase II selects a metric. A minimal record shape is shown in
Table~\ref{tab:destino_record_shape}; field names are implementation placeholders, not final Solidity declarations.

\begin{table}[H]
\centering
\footnotesize
\renewcommand{\arraystretch}{1.15}
\begin{tabularx}{\textwidth}{p{3.2cm}X X}
\toprule
\textbf{Record} & \textbf{Core fields} & \textbf{Role} \\
\midrule
\texttt{ProjectSubmission} &
\texttt{projectId}, \texttt{developerId}, \texttt{versionId}, \texttt{archiveLink}, \texttt{archiveHash},
\texttt{tags}, \texttt{storagePersistence}, \texttt{submittedAt} &
Registers the artifact/provenance claim that later metric evidence depends on. \\
\texttt{MetricRecord} &
\texttt{projectId}, \texttt{versionId}, \texttt{metricName}, \texttt{metricValue}, \texttt{evidenceHash},
\texttt{analysisPolicy}, \texttt{recordedAt} &
Stores the compact selected metric result and binds it to an auditable off-chain report. \\
\bottomrule
\end{tabularx}
\caption{Metric-agnostic record shape for the future DESTINO PoC.}
\label{tab:destino_record_shape}
\end{table}

The \texttt{analysisPolicy} field is included because the current report now makes population and sensitivity rules explicit:
human-only M3/M4 outputs, short-window flags, and bot-sensitivity summaries are not incidental cleaning details. They are part
of the evidence interpretation that must accompany any future metric record.

\subsection{Transaction and Replay Flow}
\label{sec:design_tx_flow}

The intended end-to-end flow is:
\begin{enumerate}[leftmargin=*]
  \item a developer registers signed project/archive metadata and an archive hash;
  \item an off-chain verifier retrieves the artifact and checks hash equality;
  \item the selected Phase II metric is computed using the pinned script version and declared analysis policy;
  \item a compact metric value and evidence-report hash are recorded on-chain;
  \item an event is emitted so an indexing/search component can expose the new datapoint;
  \item auditors can replay the same computation from the artifact hash, script version, window definition, and evidence report.
\end{enumerate}

This design preserves the core DESTINO requirement: trust is not based on a private metric calculation alone, but on a replayable
chain of artifact commitments, deterministic scripts, and compact public records.

\subsection{How Phase I Outputs Support the Blueprint}
\label{sec:design_to_phase1}

The current Phase I outputs are already shaped for this future integration:
\begin{itemize}[leftmargin=*]
  \item frozen projects, versions, and windows provide stable input identifiers;
  \item clean and human-only metric tables provide candidate values that can be reduced to compact records;
  \item duration flags and bot summaries document analysis policies that affect interpretation;
  \item logs and reproducibility checkpoints provide the replay trail for third-party verification;
  \item the feasibility rubric defines the gate that prevents an empirically attractive but contract-hostile metric from becoming the PoC.
\end{itemize}

This means the current report completes the DESTINO alignment at the planning and evidence level. The implementation step remains
properly deferred until Phase II identifies the metric with the best empirical and feasibility tradeoff.
